\begin{table}[htbp]
\centering
\caption{Статистическая информация для полосы $\nu_2$ молекул HD$^{32}$S и HD$^{34}$S.}
\label{tab:hds_statistics}
\begingroup
\scriptsize
\setlength{\tabcolsep}{2pt}
\renewcommand{\arraystretch}{1.05}
\begin{tabular}{@{}lccccccccc@{}}
\hline
Изотополог и источник &
Центр, см$^{-1}$ &
$J_{\max}$ &
$K_{a,\max}$ &
$N_{tr}$ &
$N_l$ &
$m_1$, \% &
$m_2$, \% &
$m_3$, \% &
$m_4$, \% \\
\hline
HD$^{32}$S, \cite{hds_ulenikov1995} & 1032.715171 & 22 & 10 & 1400 & 226 & --   & --   & --  & --  \\
HD$^{32}$S, наст. работа             & 1032.715193 & 25 & 17 & 2684 & 415 & 80.0 & 11.6 & 6.5 & 1.9 \\
HD$^{34}$S, \cite{hds_ulenikov1995} & 1031.507917 & 16 & 9  & 430  & 130 & --   & --   & --  & --  \\
HD$^{34}$S, наст. работа             & 1031.507611 & 21 & 14 & 1212 & 263 & 85.9 & 9.1  & 2.3 & 2.7 \\
\hline
\end{tabular}
\endgroup

\vspace{0.3em}
\begin{minipage}{0.98\linewidth}
\footnotesize
$N_{tr}$~--- число отнесенных переходов; $N_l$~--- число полученных энергий состояния (010).
Величины $m_i$ показывают доли уровней с остаточными отклонениями:
$m_1$~--- до $10 \times 10^{-5}$~см$^{-1}$,
$m_2$~--- от $10$ до $20 \times 10^{-5}$~см$^{-1}$,
$m_3$~--- от $20$ до $30 \times 10^{-5}$~см$^{-1}$,
$m_4$~--- более $30 \times 10^{-5}$~см$^{-1}$.
\end{minipage}
\end{table}